% ============================================================
% Section 1: Introduction
% ============================================================

\IEEEPARstart{I}{dentity} document fraud poses a growing threat to national security,
financial integrity, and border control worldwide. The widespread availability of
image-editing software has made it possible to produce convincing forgeries of
passports, national identity cards, and driver's licenses at low cost, enabling
identity theft, illegal immigration, and financial crime~\cite{TODO_idfraud_report}.
Automated forensic analysis of identity documents is therefore a critical component
of modern digital identity verification pipelines.

Detection of forged identity documents is inherently a two-level problem. At the
\emph{document level}, a system must decide whether a presented document is genuine
or has been tampered with (binary detection). At the \emph{pixel level}, it must
also identify and delineate which specific regions have been altered (forgery
localization). These two objectives are complementary: detection alone cannot
provide the evidence required for forensic audit, while localization alone is
insufficient as a security gate that must handle the entire document with a binary
decision. A practical forensic system must address both simultaneously.

\subsection*{Challenges}

The primary technical difficulty in joint detection and localization is that the
two objectives are in \emph{structural tension}: the classification head benefits
from global image features that aggregate evidence across the entire document,
while the segmentation head requires spatially precise, pixel-level representations
to delineate tampered regions. Prior work has typically addressed this tension by
either (i)~treating the two tasks independently with separate models, which loses
the opportunity for mutual regularization~\cite{trufor2023}, or (ii)~collapsing
both losses into a single scalar objective via a weighted sum, which requires
manually tuning the weight $\lambda$ before any training
begins~\cite{caruana1997multitask, kendall2018multi}. The choice of $\lambda$ is
arbitrary, dataset-dependent, and cannot be justified without exhaustive
preliminary experiments.

A second challenge is the risk of over-fitting the model-selection criterion to the
training distribution. Standard hyperparameter optimization (HPO) maximizes a
single validation metric, ignoring potential degradation on the complementary
objective. When the HPO criterion is PR-AUC, the model may sacrifice segmentation
quality to maximize detection confidence; when it is Dice coefficient, the
classification head may be under-trained. Neither outcome is acceptable for a
forensic system that must meet simultaneous performance requirements on both tracks.

\subsection*{Contributions}

This paper introduces \textbf{DocVerify}, a multi-task forensic document analysis
system that addresses both challenges through a principled multi-objective
hyperparameter optimization framework. Our contributions are as follows:

\begin{enumerate}
  \item \textbf{Multi-objective Pareto HPO for multitask forensics.}
    We propose a hyperparameter search framework in which both PR-AUC (detection)
    and Dice coefficient (localization) are treated as simultaneous optimization
    objectives. The optimal configuration is selected from the Pareto front of
    non-dominated solutions by minimizing the Euclidean distance to the ideal
    point~$(1,~1)$, a criterion that assigns equal weight to both tasks and
    requires no manual tuning of task trade-off weights.

  \item \textbf{DocVerify architecture.}
    We design an end-to-end multi-task network comprising a lightweight
    convolutional encoder (six stages, 8--256 channels) coupled to a U-Net decoder
    with skip connections for pixel-level mask prediction and a dedicated
    classification head for document-level detection. The architecture is trained
    jointly with a combined BCE+Dice segmentation loss and a BCE classification
    loss, with the loss weight $\lambda_{\text{mask}}$ included in the HPO search
    space.

  \item \textbf{Rigorous nested cross-validation protocol.}
    We evaluate DocVerify under a 10-fold outer $\times$ 5-fold inner nested
    cross-validation scheme, running 50 Optuna MOTPE trials per outer fold
    (500 model evaluations in total) entirely on the training data. The final
    performance is estimated on a blind holdout using 30 independent random seeds
    per variant, and statistical comparisons use the Wilcoxon signed-rank test with
    Holm--Bonferroni correction.

  \item \textbf{Competitive results on the DeepID 2025 Challenge benchmark.}
    Training exclusively on the FantasyID dataset~\cite{fantasyid2024} and without
    any external pre-training, DocVerify achieves a detection F1@0.5 of
    $0.969 \pm 0.014$ (Track~1) and a localization F1 of $0.807 \pm 0.096$
    (Track~2), surpassing all FantasyID-only baselines and exceeding the challenge
    winner on Track~2.
\end{enumerate}

The remainder of the paper is organized as follows. Section~\ref{sec:related}
reviews related work on document forgery detection, multi-task learning, and HPO.
Section~\ref{sec:method} describes the DocVerify architecture, loss functions, and
optimization framework. Section~\ref{sec:experiments} presents experimental results
including ablation studies and comparison with challenge participants.
Section~\ref{sec:discussion} discusses the findings, limitations, and implications.
Section~\ref{sec:conclusion} concludes the paper.
